%package list
\documentclass{article}
\usepackage[top=3cm, bottom=3cm, outer=3cm, inner=3cm]{geometry}
\usepackage{graphicx}
\usepackage{url}
\usepackage{hyperref}
\usepackage{array}
\newcolumntype{x}[1]{>{\centering\arraybackslash\hspace{0pt}}p{#1}}
\usepackage{natbib}
\usepackage{pdfpages}
\usepackage{multirow}
\usepackage[T1]{fontenc}
\usepackage{imakeidx}
\usepackage{wrapfig}
\usepackage{caption}
\usepackage{float}

\usepackage{hyperref}
\hypersetup{
    colorlinks=true,
    linkcolor=black,
    filecolor=magenta,
    urlcolor=blue,
    pdftitle={Pruebas StringUtils JUnit Mockito},
    pdfpagemode=FullScreen,
}

\urlstyle{same}

\usepackage[normalem]{ulem}
\useunder{\uline}{\ul}{}

\usepackage{minted}

\usepackage{listings}
\usepackage{color, colortbl}
\definecolor{dkgreen}{rgb}{0,0.6,0}
\definecolor{gray}{rgb}{0.5,0.5,0.5}
\definecolor{mauve}{rgb}{0.58,0,0.82}
\definecolor{codebackground}{rgb}{0.95, 0.95, 0.92}
\definecolor{tablebackground}{rgb}{0.0, 0.45, 0.63}

\lstset{
    frame=tb,
    language=bash,
    aboveskip=3mm,
    belowskip=3mm,
    showstringspaces=false,
    columns=flexible,
    basicstyle={\small\ttfamily},
    numbers=none,
    numberstyle=\tiny\color{gray},
    keywordstyle=\color{blue},
    commentstyle=\color{dkgreen},
    stringstyle=\color{mauve},
    breaklines=true,
    breakatwhitespace=true,
    tabsize=3,
    backgroundcolor=\color{codebackground},
}

%%%%%%%%%%%%%%%%%%%%%%%%%%%%%%%%%%%%%%%%%%%%%%%%%%%%%%%%%%%%%%%%%%%%%%%%%%%%
\newcommand{\csemail}{pramirezsa@ulasalle.edu.pe}
\newcommand{\csdocente}{MSc. Luis Pati\~no Herrera}
\newcommand{\cscurso}{Pruebas de Software}
\newcommand{\csuniversidad}{Universidad La Salle}
\newcommand{\csescuela}{Escuela Profesional de Ingenier\'ia de Software}
\newcommand{\cspracnr}{05}
\newcommand{\cstema}{Pruebas unitarias de StringUtils con JUnit y Mockito}
%%%%%%%%%%%%%%%%%%%%%%%%%%%%%%%%%%%%%%%%%%%%%%%%%%%%%%%%%%%%%%%%%%%%%%%%%%%%

\usepackage[english,spanish]{babel}
\usepackage[utf8]{inputenc}
\AtBeginDocument{\selectlanguage{spanish}}
\renewcommand{\figurename}{Figura}
\renewcommand{\refname}{Referencias}
\renewcommand{\tablename}{Tabla}

\usepackage{fancyhdr}
\pagestyle{fancy}
\fancyhf{}
\setlength{\headheight}{30pt}
\renewcommand{\headrulewidth}{1pt}
\renewcommand{\footrulewidth}{1pt}

\fancyhead[L]{\raisebox{-0.2\height}{\includegraphics[width=3cm]{logo_ulasalle (1).png}}}
\fancyhead[R]{\fontsize{7}{7}\selectfont \csuniversidad \\ \csescuela \\ \textbf{\cscurso}}

\fancyfoot[C]{Pruebas de Software}
\fancyfoot[R]{P\'agina \thepage}

\begin{document}

\vspace*{10px}

\begin{center}
    \fontsize{17}{17} \textbf{ Pr\'actica \cspracnr}
\end{center}

\renewcommand{\arraystretch}{1.5}

\begin{table}[h]
\begin{tabular}{|x{4.7cm}|x{4.8cm}|x{4.8cm}|}
\hline
\textbf{DOCENTE} & \textbf{CARRERA} & \textbf{CURSO} \\
\hline
\csdocente & \csescuela & \cscurso \\
\hline
\end{tabular}
\end{table}

\begin{table}[h]
\begin{tabular}{|x{4.7cm}|x{4.8cm}|x{4.8cm}|}
\hline
\textbf{GRUPO} & \textbf{TEMA} & \textbf{DURACI\'ON} \\
\hline
- & \cstema & - \\
\hline
\end{tabular}
\end{table}

\renewcommand{\arraystretch}{1}

\section*{Integrantes}
\begin{itemize}
    \item Patrick Andres Ramirez Santos
\end{itemize}

\section*{Referencia}
Documentaci\'on: \url{https://commons.apache.org/proper/commons-lang/apidocs/org/apache/commons/lang3/StringUtils.html}

\tableofcontents
\newpage

%% ============================================================
\section{Objetivo}
Implementar 20 casos de prueba para la clase \texttt{StringUtils} de Apache Commons Lang3: 10 con JUnit 5 puro y 10 con Mockito (\texttt{MockedStatic}, ya que todos los m\'etodos de \texttt{StringUtils} son est\'aticos).

%% ============================================================
\section{Configuraci\'on del proyecto}

\subsection{Estructura}
\begin{minted}{text}
Parcial/
|-- pom.xml
`-- src/
    |-- main/java/org/example/Main.java
    `-- test/java/org/example/
        |-- StringUtilsJUnitTest.java
        `-- StringUtilsMockitoTest.java
\end{minted}

\subsection{Dependencias en \texttt{pom.xml}}
\begin{minted}{xml}
<properties>
    <maven.compiler.source>21</maven.compiler.source>
    <maven.compiler.target>21</maven.compiler.target>
    <junit.version>5.11.4</junit.version>
    <mockito.version>5.14.2</mockito.version>
</properties>

<dependencies>
    <dependency>
        <groupId>org.apache.commons</groupId>
        <artifactId>commons-lang3</artifactId>
        <version>3.17.0</version>
    </dependency>
    <dependency>
        <groupId>org.junit.jupiter</groupId>
        <artifactId>junit-jupiter</artifactId>
        <version>${junit.version}</version>
        <scope>test</scope>
    </dependency>
    <dependency>
        <groupId>org.mockito</groupId>
        <artifactId>mockito-core</artifactId>
        <version>${mockito.version}</version>
        <scope>test</scope>
    </dependency>
    <dependency>
        <groupId>org.mockito</groupId>
        <artifactId>mockito-junit-jupiter</artifactId>
        <version>${mockito.version}</version>
        <scope>test</scope>
    </dependency>
</dependencies>

<build>
    <plugins>
        <plugin>
            <groupId>org.apache.maven.plugins</groupId>
            <artifactId>maven-surefire-plugin</artifactId>
            <version>3.5.2</version>
        </plugin>
    </plugins>
</build>
\end{minted}

%% ============================================================
\newpage
\section{Casos de prueba con JUnit 5}

Archivo: \texttt{src/test/java/org/example/StringUtilsJUnitTest.java}

\begin{table}[H]
\centering
\begin{tabular}{|x{1cm}|x{4cm}|x{8cm}|}
\hline
\textbf{\#} & \textbf{M\'etodo} & \textbf{Qu\'e valida} \\
\hline
1 & \texttt{isEmpty} & \texttt{null} y \texttt{""} devuelven \texttt{true}; espacios y texto, \texttt{false}. \\
\hline
2 & \texttt{isBlank} & \texttt{null}, \texttt{""} y solo espacios devuelven \texttt{true}. \\
\hline
3 & \texttt{capitalize} & Pone en may\'uscula la primera letra; \texttt{null}-safe. \\
\hline
4 & \texttt{reverse} & Invierte la cadena; \texttt{null}-safe. \\
\hline
5 & \texttt{join} & Une un arreglo con separador \texttt{String} o \texttt{char}. \\
\hline
6 & \texttt{repeat} & Repite N veces; N $\leq$ 0 devuelve \texttt{""}. \\
\hline
7 & \texttt{contains} & Verifica subcadena; \texttt{null} devuelve \texttt{false}. \\
\hline
8 & \texttt{substring} & Soporta \'indices negativos y fuera de rango. \\
\hline
9 & \texttt{swapCase} & Invierte mayus/minus de cada letra. \\
\hline
10 & \texttt{countMatches} & Cuenta ocurrencias de subcadena. \\
\hline
\end{tabular}
\caption{Casos de prueba con JUnit puro}
\end{table}

\subsection{Ejemplo}
\begin{minted}{java}
@Test
@DisplayName("isBlank: true para null, \"\" y solo espacios")
void testIsBlank() {
    assertTrue(StringUtils.isBlank(null));
    assertTrue(StringUtils.isBlank(""));
    assertTrue(StringUtils.isBlank("   "));
    assertFalse(StringUtils.isBlank("a"));
    assertFalse(StringUtils.isBlank("  a  "));
}
\end{minted}

\begin{minted}{bash}
mvn test -Dtest=StringUtilsJUnitTest
\end{minted}

Se ejecuta solamente la clase \texttt{StringUtilsJUnitTest}. Los 10 casos pasan sin fallos ni errores.

\begin{figure}[H]
\centering
% \includegraphics[width=0.9\linewidth]{captura_junit.png}
\fbox{\parbox{0.9\linewidth}{\centering \vspace{2cm} [Captura: ejecuci\'on de los 10 tests JUnit] \vspace{2cm}}}
\caption{Resultado de los tests con JUnit puro}
\end{figure}

%% ============================================================
\newpage
\section{Casos de prueba con Mockito}

Archivo: \texttt{src/test/java/org/example/StringUtilsMockitoTest.java}

Como \texttt{StringUtils} solo expone m\'etodos est\'aticos, se usa \texttt{Mockito.mockStatic(...)} dentro de un \texttt{try-with-resources} para acotar el alcance del mock.

\begin{table}[H]
\centering
\begin{tabular}{|x{1cm}|x{4cm}|x{8cm}|}
\hline
\textbf{\#} & \textbf{M\'etodo mockeado} & \textbf{Qu\'e valida} \\
\hline
11 & \texttt{isEmpty} & Stub devuelve \texttt{true} para entrada no vac\'ia. \\
\hline
12 & \texttt{isBlank} & Stub devuelve \texttt{false} para \texttt{""}. \\
\hline
13 & \texttt{capitalize} & Stub con \texttt{anyString()} devuelve valor fijo. \\
\hline
14 & \texttt{reverse} & Stub mantiene la misma cadena. \\
\hline
15 & \texttt{join} & Stub controla el resultado completo. \\
\hline
16 & \texttt{repeat} & Stub ignora el conteo real. \\
\hline
17 & \texttt{contains} & Stub fuerza \texttt{true} con verificaci\'on por \texttt{eq()}. \\
\hline
18 & \texttt{substring} & Stub devuelve un fragmento simulado. \\
\hline
19 & \texttt{swapCase} & Stub con \texttt{anyString()} y \texttt{verify} con \texttt{times(2)}. \\
\hline
20 & \texttt{countMatches} & Stub devuelve un valor fijo y verifica la llamada. \\
\hline
\end{tabular}
\caption{Casos de prueba con Mockito (\texttt{MockedStatic})}
\end{table}

\subsection{Ejemplo}
\begin{minted}{java}
@Test
@DisplayName("Mockito: stub de capitalize devuelve un valor fijo")
void testMockedCapitalize() {
    try (MockedStatic<StringUtils> mocked = mockStatic(StringUtils.class)) {
        mocked.when(() -> StringUtils.capitalize(anyString()))
              .thenReturn("MOCKED");

        assertEquals("MOCKED", StringUtils.capitalize("cualquiera"));
        assertEquals("MOCKED", StringUtils.capitalize("otra"));
        mocked.verify(() -> StringUtils.capitalize(anyString()), times(2));
    }
}
\end{minted}

\begin{minted}{bash}
mvn test -Dtest=StringUtilsMockitoTest
\end{minted}

Se ejecuta solamente la clase \texttt{StringUtilsMockitoTest}. Los 10 casos pasan; aparece el warning esperado del \textit{agent} de ByteBuddy al cargar din\'amicamente el \textit{inline-mock-maker} en JDK 21+.

\begin{figure}[H]
\centering
% \includegraphics[width=0.9\linewidth]{captura_mockito.png}
\fbox{\parbox{0.9\linewidth}{\centering \vspace{2cm} [Captura: ejecuci\'on de los 10 tests Mockito] \vspace{2cm}}}
\caption{Resultado de los tests con Mockito}
\end{figure}

%% ============================================================
\newpage
\section{Ejecuci\'on}

\subsection{Comando}
\begin{minted}{bash}
mvn test
\end{minted}

\subsection{Salida}
\begin{minted}{text}
[INFO] Running org.example.StringUtilsJUnitTest
[INFO] Tests run: 10, Failures: 0, Errors: 0, Skipped: 0
[INFO] Running org.example.StringUtilsMockitoTest
[INFO] Tests run: 10, Failures: 0, Errors: 0, Skipped: 0
[INFO] Results:
[INFO] Tests run: 20, Failures: 0, Errors: 0, Skipped: 0
[INFO] BUILD SUCCESS
\end{minted}

\begin{minted}{bash}
mvn clean test
\end{minted}

Compilaci\'on limpia y ejecuci\'on de ambas clases en una sola corrida. Total: 20 tests, 0 fallos, 0 errores, \texttt{BUILD SUCCESS}.

\begin{figure}[H]
\centering
% \includegraphics[width=0.9\linewidth]{captura_build.png}
\fbox{\parbox{0.9\linewidth}{\centering \vspace{2cm} [Captura: BUILD SUCCESS con 20 tests] \vspace{2cm}}}
\caption{Salida final de \texttt{mvn clean test}}
\end{figure}

%% ============================================================
\newpage
\section{Conclusiones}
\begin{itemize}
    \item JUnit 5 es suficiente para validar el comportamiento real de m\'etodos est\'aticos puros como los de \texttt{StringUtils}.
    \item Mockito agrega valor cuando se necesita aislar dependencias; con \texttt{MockedStatic} es posible interceptar m\'etodos est\'aticos dentro de un \textit{scope} controlado.
    \item Los 20 casos pasan en una sola corrida con \texttt{mvn test}.
\end{itemize}

\end{document}
